% !TEX program = xelatex
% Standalone viewer for the paired-Brier CI table.
% The table float itself lives in brier_ci_table_body.tex, generated by
% make_brier_ci_table.R, and is \input-able directly into the thesis.
\documentclass[12pt]{article}

\usepackage{fontspec}
\usepackage{xeCJK}
\setCJKmainfont[AutoFakeSlant=.2, AutoFakeBold=3]{AR PL KaitiM Big5}

\usepackage[margin=1in]{geometry}
\usepackage{amsmath, amssymb}
\usepackage{booktabs}
\usepackage{array}
\usepackage{hyperref}
\hypersetup{colorlinks=true, linkcolor=blue!45!black, urlcolor=blue!45!black}

\title{\textbf{論文用表：配對 Brier 差與站層 bootstrap 信賴區間}}
\author{取代 thesis Table \ref{tab:ozone-brier-ci-ref} 的候選}
\date{2026 年 7 月}

\begin{document}
\maketitle

本檔示範一張可直接取代論文「Top-4 名次表」的表格。原名次表把
Rank 1--4 都印成相同的 $0.030$，讀者無從判斷延伸是否有效；下表改列每個延伸
相對其\textbf{配對基準}的 Brier 差 $\Delta$ 與 $95\%$ 信賴區間，一眼即可看出
\emph{所有區間（除最稀疏一格外）都涵蓋 $0$}。

\bigskip

% ---- 這裡 \input 由 R 產生的表格 float；論文中把下一行搬進 mythesis.tex 即可 ----
% Auto-generated by make_brier_ci_table.R -- do not edit by hand.
\begin{table}[htbp]
\centering
\caption{Paired held-out Brier differences for each extension against its
matched baseline, $\Delta = \mathrm{BS}_{\text{ext}} - \mathrm{BS}_{\text{base}}$
(negative favours the extension), in units of $10^{-3}$, with $95\%$
site-clustered bootstrap confidence intervals in brackets ($B=2000$).
AR(2) vs.\ AR(1): skew-$t$, $K{=}7$, $T{=}50$; MRTS vs.\ none: skew-$t$,
$K{=}1$, $T{=}0$. $n_{\mathrm{exc}}$ is the number of exceedance events
behind the score. Every interval covers $0$ except the starred one
($q{=}0.995$, $200/200$), which sits at the sparsest level and does not
survive multiplicity across the $15$ comparisons.}
\label{tab:ozone-brier-ci}
\small
\begin{tabular}{@{}ll r@{\ }l r@{\ }l@{}}
\toprule
& & \multicolumn{2}{c}{$200/200$ split} & \multicolumn{2}{c}{$300/100$ split} \\
\cmidrule(lr){3-4}\cmidrule(lr){5-6}
Comparison & $L$ quantile & $\Delta\times10^{3}$ [95\% CI] & $n_{\mathrm{exc}}$ & $\Delta\times10^{3}$ [95\% CI] & $n_{\mathrm{exc}}$ \\
\midrule
AR(2) vs.\ AR(1) & $0.900$ & $+0.10$ {\scriptsize $[-0.35,\,+0.54]$} & 630 & $-0.03$ {\scriptsize $[-0.53,\,+0.45]$} & 241 \\
 & $0.950$ & $-0.06$ {\scriptsize $[-0.42,\,+0.29]$} & 319 & $-0.03$ {\scriptsize $[-0.33,\,+0.32]$} & 92 \\
 & $0.980$ & $-0.09$ {\scriptsize $[-0.25,\,+0.09]$} & 129 & $-0.07$ {\scriptsize $[-0.28,\,+0.13]$} & 28 \\
 & $0.990$ & $+0.02$ {\scriptsize $[-0.09,\,+0.18]$} & 56 & $+0.01$ {\scriptsize $[-0.09,\,+0.10]$} & 9 \\
 & $0.995$ & $-0.03$ {\scriptsize $[-0.08,\,+0.02]$} & 27 & $+0.06$ {\scriptsize $[-0.02,\,+0.16]$} & 2 \\
\addlinespace
MRTS vs.\ none & $0.900$ & $-0.07$ {\scriptsize $[-0.48,\,+0.35]$} & 630 & \multicolumn{1}{c}{--} & -- \\
 & $0.950$ & $-0.08$ {\scriptsize $[-0.34,\,+0.17]$} & 319 & \multicolumn{1}{c}{--} & -- \\
 & $0.980$ & $-0.11$ {\scriptsize $[-0.25,\,+0.00]$} & 129 & \multicolumn{1}{c}{--} & -- \\
 & $0.990$ & $-0.05$ {\scriptsize $[-0.12,\,+0.01]$} & 56 & \multicolumn{1}{c}{--} & -- \\
 & $0.995$ & $-0.03^{*}$ {\scriptsize $[-0.06,\,-0.00]$} & 27 & \multicolumn{1}{c}{--} & -- \\
\bottomrule
\end{tabular}
\end{table}


\bigskip
\noindent\textbf{讀法。}$\Delta<0$ 表示延伸較佳。$\Delta$ 與 CI 以 $10^{-3}$ 為
單位；$\Delta$ 約 $10^{-1}\times10^{-3}$，而 CI 半寬約 $3\text{--}5\times
10^{-1}\times10^{-3}$，差異比不確定性小一個數量級。兩個 split 結論一致：AR(2)
與 MRTS 相對配對基準\textbf{測不出差異}。唯一星號（$q{=}0.995$, $200/200$，
$n_{\mathrm{exc}}{=}27$）出現在最稀疏 level，且在 $15$ 個比較的多重性下不足採信
——$300/100$ 的同一 level 只有 $2$ 個超標事件，正說明此處雜訊之大。

\medskip
\noindent\textbf{嵌入論文步驟。}（1）把 \verb|brier_ci_table_body.tex| 複製到論文
目錄；（2）在 \verb|mythesis.tex| 以 \verb|% Auto-generated by make_brier_ci_table.R -- do not edit by hand.
\begin{table}[htbp]
\centering
\caption{Paired held-out Brier differences for each extension against its
matched baseline, $\Delta = \mathrm{BS}_{\text{ext}} - \mathrm{BS}_{\text{base}}$
(negative favours the extension), in units of $10^{-3}$, with $95\%$
site-clustered bootstrap confidence intervals in brackets ($B=2000$).
AR(2) vs.\ AR(1): skew-$t$, $K{=}7$, $T{=}50$; MRTS vs.\ none: skew-$t$,
$K{=}1$, $T{=}0$. $n_{\mathrm{exc}}$ is the number of exceedance events
behind the score. Every interval covers $0$ except the starred one
($q{=}0.995$, $200/200$), which sits at the sparsest level and does not
survive multiplicity across the $15$ comparisons.}
\label{tab:ozone-brier-ci}
\small
\begin{tabular}{@{}ll r@{\ }l r@{\ }l@{}}
\toprule
& & \multicolumn{2}{c}{$200/200$ split} & \multicolumn{2}{c}{$300/100$ split} \\
\cmidrule(lr){3-4}\cmidrule(lr){5-6}
Comparison & $L$ quantile & $\Delta\times10^{3}$ [95\% CI] & $n_{\mathrm{exc}}$ & $\Delta\times10^{3}$ [95\% CI] & $n_{\mathrm{exc}}$ \\
\midrule
AR(2) vs.\ AR(1) & $0.900$ & $+0.10$ {\scriptsize $[-0.35,\,+0.54]$} & 630 & $-0.03$ {\scriptsize $[-0.53,\,+0.45]$} & 241 \\
 & $0.950$ & $-0.06$ {\scriptsize $[-0.42,\,+0.29]$} & 319 & $-0.03$ {\scriptsize $[-0.33,\,+0.32]$} & 92 \\
 & $0.980$ & $-0.09$ {\scriptsize $[-0.25,\,+0.09]$} & 129 & $-0.07$ {\scriptsize $[-0.28,\,+0.13]$} & 28 \\
 & $0.990$ & $+0.02$ {\scriptsize $[-0.09,\,+0.18]$} & 56 & $+0.01$ {\scriptsize $[-0.09,\,+0.10]$} & 9 \\
 & $0.995$ & $-0.03$ {\scriptsize $[-0.08,\,+0.02]$} & 27 & $+0.06$ {\scriptsize $[-0.02,\,+0.16]$} & 2 \\
\addlinespace
MRTS vs.\ none & $0.900$ & $-0.07$ {\scriptsize $[-0.48,\,+0.35]$} & 630 & \multicolumn{1}{c}{--} & -- \\
 & $0.950$ & $-0.08$ {\scriptsize $[-0.34,\,+0.17]$} & 319 & \multicolumn{1}{c}{--} & -- \\
 & $0.980$ & $-0.11$ {\scriptsize $[-0.25,\,+0.00]$} & 129 & \multicolumn{1}{c}{--} & -- \\
 & $0.990$ & $-0.05$ {\scriptsize $[-0.12,\,+0.01]$} & 56 & \multicolumn{1}{c}{--} & -- \\
 & $0.995$ & $-0.03^{*}$ {\scriptsize $[-0.06,\,-0.00]$} & 27 & \multicolumn{1}{c}{--} & -- \\
\bottomrule
\end{tabular}
\end{table}
| 取代舊的
Table~7；（3）確認論文前導已有 \verb|booktabs|；（4）把
\verb|\label{tab:ozone-brier-ci}| 接到既有引用。表格由
\verb|make_brier_ci_table.R| 自資料重生，數字不需手動維護。

\vfill
\footnotesize
\noindent 參考標籤（本檔自身）：\label{tab:ozone-brier-ci-ref}
資料來源 \verb|brier_bootstrap_pairs_{200-200,300-100}.csv|，
由 \verb|code/analysis/ozone/US-all-auto/brier_bootstrap.R| 產生。

\end{document}
